% Options for packages loaded elsewhere
\PassOptionsToPackage{unicode}{hyperref}
\PassOptionsToPackage{hyphens}{url}
\documentclass[
]{ltjsarticle}
\usepackage{xcolor}
\usepackage{amsmath,amssymb}
\setcounter{secnumdepth}{-\maxdimen} % remove section numbering
\usepackage{iftex}
\ifPDFTeX
  \usepackage[T1]{fontenc}
  \usepackage[utf8]{inputenc}
  \usepackage{textcomp} % provide euro and other symbols
\else % if luatex or xetex
  \usepackage{unicode-math} % this also loads fontspec
  \defaultfontfeatures{Scale=MatchLowercase}
  \defaultfontfeatures[\rmfamily]{Ligatures=TeX,Scale=1}
\fi
\usepackage{lmodern}
\ifPDFTeX\else
  % xetex/luatex font selection
\fi
% Use upquote if available, for straight quotes in verbatim environments
\IfFileExists{upquote.sty}{\usepackage{upquote}}{}
\IfFileExists{microtype.sty}{% use microtype if available
  \usepackage[]{microtype}
  \UseMicrotypeSet[protrusion]{basicmath} % disable protrusion for tt fonts
}{}
\makeatletter
\@ifundefined{KOMAClassName}{% if non-KOMA class
  \IfFileExists{parskip.sty}{%
    \usepackage{parskip}
  }{% else
    \setlength{\parindent}{0pt}
    \setlength{\parskip}{6pt plus 2pt minus 1pt}}
}{% if KOMA class
  \KOMAoptions{parskip=half}}
\makeatother
\usepackage{longtable,booktabs,array}
\newcounter{none} % for unnumbered tables
\usepackage{calc} % for calculating minipage widths
% Correct order of tables after \paragraph or \subparagraph
\usepackage{etoolbox}
\makeatletter
\patchcmd\longtable{\par}{\if@noskipsec\mbox{}\fi\par}{}{}
\makeatother
% Allow footnotes in longtable head/foot
\IfFileExists{footnotehyper.sty}{\usepackage{footnotehyper}}{\usepackage{footnote}}
\makesavenoteenv{longtable}
\setlength{\emergencystretch}{3em} % prevent overfull lines
\providecommand{\tightlist}{%
  \setlength{\itemsep}{0pt}\setlength{\parskip}{0pt}}
\usepackage{newunicodechar}
\newunicodechar{→}{\ensuremath{\rightarrow}}
\newunicodechar{⟨}{\ensuremath{\langle}}
\newunicodechar{⟩}{\ensuremath{\rangle}}
\newunicodechar{⇒}{\ensuremath{\Rightarrow}}
\newunicodechar{⟹}{\ensuremath{\Longrightarrow}}
\newunicodechar{×}{\ensuremath{\times}}
\newunicodechar{≈}{\ensuremath{\approx}}
\newunicodechar{∼}{\ensuremath{\sim}}
\newunicodechar{≤}{\ensuremath{\le}}
\newunicodechar{≥}{\ensuremath{\ge}}
\newunicodechar{≠}{\ensuremath{\ne}}
\newunicodechar{∫}{\ensuremath{\int}}
\newunicodechar{−}{\ensuremath{-}}
\newunicodechar{✗}{\ensuremath{\times}}
\newunicodechar{✓}{\ensuremath{\checkmark}}
\usepackage{bookmark}
\IfFileExists{xurl.sty}{\usepackage{xurl}}{} % add URL line breaks if available
\urlstyle{same}
\hypersetup{
  hidelinks,
  pdfcreator={LaTeX via pandoc}}

\author{}
\date{}

\begin{document}

\section{観測者--系ビートの等分布によるボルン統計の創発 ―
局在核モデル上の予想}\label{ux89b3ux6e2cux8005ux7cfbux30d3ux30fcux30c8ux306eux7b49ux5206ux5e03ux306bux3088ux308bux30dcux30ebux30f3ux7d71ux8a08ux306eux5275ux767a-ux5c40ux5728ux6838ux30e2ux30c7ux30ebux4e0aux306eux4e88ux60f3}

\textbf{木原 範昭}
（査読対話ノート由来。\textbf{本稿は証明論文でも観察論文でもなく、予想（予言）論文である。}
確立物理の主張ではなく、明示された機構の候補と、その検証可能な帰結の提示である。）

バージョン v0.2（2026-06-27） DOI (Version): 10.5281/zenodo.20967082 DOI
(Concept): 10.5281/zenodo.20967081 Zenodo:
https://zenodo.org/records/20967082

\textbf{v0.2 改訂注}（第二査読＝Claude Code 反映）：(1)
\textbf{両ケースの厳密導出を明示}――局在入力 \(\psi^{\rm loc}\)
と広がり入力 \(\cos\varphi\)
がともに同一の再生核から出て、出力差は入力の差のみ（§3、§5(B)
で実畳み込みにより機械精度確認）;(2) \textbf{§3 (I)
の役割分担を書き直し}――等分布（Weyl）＝一様測度の供給（偏りを作らない）、\(|\psi|^2\)
の形＝再生核の決定論的読み、強度→クリック＝橋(§3
II)、の三分離を明示し「Weyl だけで二乗の偏りが出る」という混同を除去;(3)
\textbf{反証バッテリー（§5
E）}――広がり入力が尖らない・局在入力が尖る・Weyl
単独は平坦・両ケースは同一恒等式、の四条件をすべてクリア（整合性の積極的証拠）;(4)
§5(B) を \textbf{{[}2{]}
の再掲から実畳み込みの自己完結検証へ格上げ}（タウトロジー解消・必須化）。\(S_N\)
を実際に数値畳み込みし、cos は全 \(N\) 厳密・多モード（1,3,5）は
\(N\ge5\) 厳密／\(N<5\) 取りこぼしというバンド条件の \(N\)
依存も明示;(5) {[}2{]} の確定 DOI
追記;(6)「証明済みの核」が等分布の仮定（§8
O1）に条件づくこと、等分布機構が離散測定イベント列を前提とするため B1 が
B2・B3 の論理的上流にある点（§2・§8 O2）を明示;(7) {[}7{]}
の単名「Lynnx」が arXiv 記載どおりと確認;(8) \textbf{§5(D)
の橋MCの地位を明示}――D2 のモンテカルロは橋(§3
II)を\textbf{仮定して実装した整合性デモ}であって橋の導出・検証ではないことを本文・コードコメントに明記し、核（Weyl×再生核＝証明済み）と橋（強度→確率＝予想）の地位分離を維持。\textbf{主張・結論・数式は不変}（厳密化と正直化のみ）。検証コード
\texttt{born\_beat\_conjecture.py}
も両ケース畳み込み・バンド条件・役割分離(D)・反証(E)を追加。

\begin{center}\rule{0.5\linewidth}{0.5pt}\end{center}

\subsection{要旨}\label{ux8981ux65e8}

先行する局在核論文 {[}2{]} は、半波長位相区間
\(\varphi\in[-\pi/2,\pi/2]\) 上の一定振幅奇数倍音和（孤立ピーク波
\(S_N\)）の\textbf{再生核性}から、観測分布の\textbf{形}がベース波の
\(|\psi_{\rm base}|^2\)
に厳密一致することを示した。ただしそこでは「ランダム性の存在」――なぜ決定論的に計算される干渉強度プロファイルが、試行ごとに一点へ局在し、繰り返しではじめて
\(|\psi_{\rm base}|^2\)
の頻度として現れるのか――を未解決の公準として残した。

本稿はこの\textbf{ランダム性の起源}に対する機構の候補を\textbf{予想}として提示する。中心となる観察は、局在核論文の走査変数
\(\varphi_0\)
が\textbf{自由変数ではなく観測器の位相}だという点である。観測器はそれ自身が有限振動数
\(\psi\)
をもつ物理対象であり、系の外に絶対位相参照は存在しない（主観空間の前提）。ゆえに測定列は系の振動数
\(\nu\) と観測器の振動数 \(\psi\)
の\textbf{相対位相}（ビート／ヘテロダイン）\(\,(\nu-\psi)t\,\)
を走査する。相対位相比が無理数のとき、この相対位相は区間上で\textbf{等分布する}（Weyl
の一様分布定理）。

本稿の論理は二段からなる。\textbf{(I)
条件付きで証明済みの核}：再生核性（{[}2{]}）により各観測器位相
\(\varphi_0\) での二乗オーバラップは厳密に
\((\pi/2)^2|\psi_{\rm base}(\varphi_0)|^2\)
であり、ビート位相が等分布する（無理比、§8
O1）ならば、長時間で得られる\textbf{位置分解された強度プロファイル}が規格化後
\(|\psi_{\rm base}|^2\)
に一致する。ここで保証されるのは決定論的に掃引される\textbf{強度プロファイルの形}であって、単発確率ではない。この帰結は、等分布の仮定（O1）を認めれば、再生核性（証明済み）と
Weyl の定理（証明済み）から追加の仮定なく従う。\textbf{(II)
予想（橋）}：この強度プロファイルが、閾値型検出を経て\textbf{単発クリックの頻度}
\(\propto|\psi_{\rm base}|^2\) として現れる、という移行を予想する。(II)
は本稿が\textbf{導出しない}部分であり、なぜ一発ごとに単一点が立つか（離散性）はここでは未解決として残す。

この予想は孤立した着想ではなく、二つの確立した系譜の\textbf{交点}に位置する：(a)
Khrennikov の prequantum classical statistical field
theory（PCSFT）と閾値検出によるボルン則の創発
{[}3--6{]}（物理的描像）、(b) Poincaré に始まり Feintzeig
らが量子へ適用した\textbf{任意関数の方法}（method of arbitrary
functions）{[}7--10{]}（数学的機構＝本稿の等分布）。本稿固有の寄与は、両系譜が抽象的に置く「等分布する高速確率変数」を、観測器--系の\textbf{ヘテロダイン・ビート}として物理的に同定し、かつ再生核性により強度プロファイルの形を\textbf{厳密}に固定する点にある。さらに、有限
\(N\)（離散格子）での再生のずれ \(O(1/N)\)
が\textbf{ボルン則からの測定可能な逸脱}を与えるという、反証可能な予言を提示する。

\begin{center}\rule{0.5\linewidth}{0.5pt}\end{center}

\subsection{1.
背景と本稿の位置づけ}\label{ux80ccux666fux3068ux672cux7a3fux306eux4f4dux7f6eux3065ux3051}

局在核論文 {[}2{]} は、ボルン則 \(P\propto|\psi|^2\)
の二乗構造のうち「形」を再生核性から導き、残す公準を二点に絞った：(a)
振幅→確率の\textbf{二乗則}（指数2）、(c)
\textbf{ランダム性の存在}。さらに測定を局在核の畳み込みで表す\textbf{測定モデル}(b)
も前提として残る。{[}2{]} が厳密に確立したのは次である：帯域制限ベース波
\(\psi_{\rm base}=\sum_{|k|\le N}c_k e^{ik\varphi}\)
に対し、シフト局在核 \(S_N\) との畳み込みが

\[
(S_N*\psi_{\rm base})(\varphi_0)=\frac{\pi}{2}\,\psi_{\rm base}(\varphi_0)
\quad\Longrightarrow\quad
\big|(S_N*\psi_{\rm base})(\varphi_0)\big|^2=\frac{\pi^2}{4}\,|\psi_{\rm base}(\varphi_0)|^2
\]

を全域・全 \(N\) で厳密に満たす（複素の場合は真のモジュラス
\(Z\bar Z\)）。

本稿は (c)
を主題とする。これは三つの公準のうち最も深く、標準理論でも未解決の測定問題そのものである。本稿はこれを\textbf{解いたとは主張せず}、機構の候補を\textbf{予想}として定式化し、その一部（強度プロファイルの再現）が証明済みの核から従うことを示し、残る部分（単発性・確率解釈）を明示的に未解決として切り出す。

\begin{center}\rule{0.5\linewidth}{0.5pt}\end{center}

\subsection{\texorpdfstring{2. 問題の精密化：\(\varphi_0\)
は誰が決めるか}{2. 問題の精密化：\textbackslash varphi\_0 は誰が決めるか}}\label{ux554fux984cux306eux7cbeux5bc6ux5316varphi_0-ux306fux8ab0ux304cux6c7aux3081ux308bux304b}

{[}2{]} では \(\varphi_0\)
は区間上を動く自由変数として全域がプロットされる。しかし\textbf{一回の観測は一つの
\(\varphi_0\) でしか起こらない}。決定論的計算が与えるのは関数
\(|\psi_{\rm base}(\cdot)|^2\)
だが、一回の測定は単一の実現値を返す。この「\textbf{関数 →
一回ごとの単一実現値}」の移行が、残された謎の全部である。これを三つの論理的に独立な橋に分解する：

\begin{itemize}
\tightlist
\item
  \textbf{B1（離散性）}：なぜ毎回ひとつの結果だけが立つのか（全プロファイルが一斉に出ない）。
\item
  \textbf{B2（どれか）}：離散だとして、なぜどの \(\varphi_0\)
  かが状態から予測不能なのか。
\item
  \textbf{B3（集団＝強度）}：なぜ試行頻度が強度プロファイルに一致するのか。
\end{itemize}

重み
\(\propto|\)重なり\(|^2\)（黄金律的な結合率）は三つの中で最も非自明でない。位置のゆらぎは
B1・B2 の現れであり、\(|\psi|^2\) の\textbf{形}の謎ではない――形は
{[}2{]} で強度として既に片付いている。本稿の予想は主に B2 と B3
に答え、B1 を未解決として残す。

ただし三つの橋は\textbf{完全には独立でない}。後述（§3--§4）の等分布機構は\textbf{離散的な測定イベント列}の上で初めて非自明になる――連続時間では
\(\varphi_0=(\nu-\psi)t\)
の掃引が無理性なしに一様化し、有理比の反例も消える――ため、B1（離散性）は
B2・B3
の機構の論理的に\textbf{上流}にある。すなわち等分布論法は離散クロックを暗黙に前提とする。この依存関係を
§8 O2 で明示する。

\begin{center}\rule{0.5\linewidth}{0.5pt}\end{center}

\subsection{\texorpdfstring{3. 予想：観測者--系ビートによる
\(\varphi_0\)
の等分布}{3. 予想：観測者--系ビートによる \textbackslash varphi\_0 の等分布}}\label{ux4e88ux60f3ux89b3ux6e2cux8005ux7cfbux30d3ux30fcux30c8ux306bux3088ux308b-varphi_0-ux306eux7b49ux5206ux5e03}

\textbf{前提（追加する唯一の物理）}。観測器は有限振動数 \(\psi\)
をもつ物理対象である。完全な絶対位相参照（系外の絶対時計）は存在しない――これは
{[}2{]}
と本稿が共有する主観空間の前提（差のみが物理的・絶対位相は非可観測）の必然であって、新たな仮定ではない。観測器が完全参照なら
\(\varphi_0\)
は固定され、ゆらぎは消えるが、そのような器は前提が禁じている。

\textbf{ビート}。系 \(\cos(\nu t)\) と観測器 \(\cos(\psi t)\) の積は

\[
\cos(\nu t)\cos(\psi t)=\tfrac12\big[\cos((\nu-\psi)t)+\cos((\nu+\psi)t)\big]
\]

で、差周波 \(\nu-\psi\)
がゆっくりした包絡（干渉縞の周期）を与える。測定タイミングごとに観測器位相が進むため、有効な走査位相は
\(\varphi_0=(\nu-\psi)t\)
となる。観測器が系と同一振動数（\(\psi=\nu\)）なら \(\varphi_0\)
は固定され、毎回同じ位相を読む――縞は出ず、単一値。観測器が系とずれる（\(\psi\ne\nu\)）と、\(\varphi_0\)
がスキャンされ、繰り返しが分布を掃く。

\textbf{単発の決定性}。重要な点として、各 \(\varphi_0\)
における一回の畳み込み読み出しは {[}2{]}
により\textbf{決定論的}である：\((S_N*\psi_{\rm base})(\varphi_0)=\tfrac{\pi}{2}\psi_{\rm base}(\varphi_0)\)
はシャープで曖昧でない。これは観測されるピーク波が実際にシャープに局在する事実（ランダムには散らばらない）と整合する。曖昧さは単発にではなく、\textbf{繰り返しにわたる
\(\varphi_0\) の分布}に宿る。

\textbf{予想（本体）}。測定が（観測器クロックにより）離散的な時刻
\(t_n\) で起こり、相対位相比 \(r=(\nu-\psi)\,\Delta t/(2\pi)\)
が無理数のとき、\(\varphi_0=(\nu-\psi)t_n\) は区間 \([-\pi/2,\pi/2]\)
上で\textbf{等分布する}（無理回転に対する Weyl
の一様分布定理）。離散測定列を前提とする点（連続時間では無理性は不要・反例も消える）は
§2 の依存関係に対応する。ゆえに長時間で各 \(\varphi_0\)
は一様に標本化され、位置分解された二乗読み出しの長時間平均が

\[
\boxed{\;\lim_{T\to\infty}\frac1T\int_0^T \big|(S_N*\psi_{\rm base})\big((\nu-\psi)t\big)\big|^2\,\delta_{\varphi_0}\,dt
\;\propto\; |\psi_{\rm base}(\varphi_0)|^2\;}
\]

すなわち\textbf{位置分解強度プロファイル}が規格化後
\(|\psi_{\rm base}|^2\) に一致する。この等式は、再生核性（§1
の厳密結果）と Weyl
等分布（古典定理）の二つから、\textbf{等分布の仮定（§8
O1）の下で追加の仮定なく従う}（§5
で数値確認）。保証されるのは強度プロファイルの\textbf{形}であり、単発確率ではない（橋は下記）。

\textbf{役割分担（混同の除去）。}
ここで三つの寄与を厳密に分ける。\textbf{等分布（Weyl）は一様測度
\(d\varphi_0\)
を供給するだけで、それ自体は二乗の偏りを作らない}――観測器位相
\(\varphi_0\)
が一様に走査されれば全位置は等頻度（定数）であって、\(|\psi|^2\)
の偏りにはならない。\(|\psi_{\rm base}(\varphi_0)|^2\)
という\textbf{位置ごとの重み（形）の出どころは、各 \(\varphi_0\)
での再生核の決定論的読み}（{[}2{]}：\((S_N*\psi_{\rm base})(\varphi_0)=\tfrac{\pi}{2}\psi_{\rm base}(\varphi_0)\)、入力
\(\psi_{\rm base}\)
をそのまま再生）であって、等分布ではない。そして\textbf{強度プロファイルを単発クリックの頻度へ変えるのは閾値検出の橋}（§3
(II)）である。整理すると:

\[
\underbrace{|\psi_{\rm base}(\varphi_0)|^2\ \text{の頻度}}_{\text{ボルン出力}}
=\underbrace{d\varphi_0\ (\text{一様})}_{\text{Weyl 等分布}}\ \times\
\underbrace{|\psi_{\rm base}(\varphi_0)|^2}_{\text{再生核の決定論的読み（形）}}\ \times\
\underbrace{(\text{強度}\to\text{クリック})}_{\text{橋・§3 (II)}}.
\]

\textbf{証明済みの核}が主張するのは中央の因子（強度プロファイルの形
\(=|\psi_{\rm base}|^2\)）に限り、左の一様測度を Weyl
が供給する。右の橋は予想である。Weyl
単独で偏りが出るかのような書き方はしない（§5 D で数値的にも、等分布した
\(\varphi_0\)
のヒストグラムは\textbf{平坦}で、強度重み付けを通して初めて \(|\psi|^2\)
が出ることを確認する）。

\textbf{両ケースの厳密導出（同一機構・入力依存のみ）。} 再生核恒等式
\((S_N*\psi_{\rm base})(\varphi_0)=\tfrac{\pi}{2}\psi_{\rm base}(\varphi_0)\)
は入力 \(\psi_{\rm base}\) に依らず成り立つので、二つの代表入力に対し:

\begin{itemize}
\tightlist
\item
  \textbf{ケース1（局在）}：\(\psi_{\rm base}^{\rm loc}=\sum_{|k|\le N}c_k e^{ik\varphi}\)（多モード重ね合わせ、中央に局在。最も局在した例は
  \(S_N\) 自身）に対し
  \[\big|(S_N*\psi_{\rm base}^{\rm loc})(\varphi_0)\big|^2=\tfrac{\pi^2}{4}\,|\psi_{\rm base}^{\rm loc}(\varphi_0)|^2\ (\text{尖った出力}).\]
\item
  \textbf{ケース2（広がり）}：\(\psi_{\rm base}^{\cos}=\cos\varphi\)（広がった単一モード）に対し
  \[\big|(S_N*\cos)(\varphi_0)\big|^2=\tfrac{\pi^2}{4}\cos^2\varphi_0\ (\text{広がった出力}).\]
\end{itemize}

両者は\textbf{同一の再生核}から出て、出力の差は入力 \(\psi_{\rm base}\)
の差のみである。\textbf{広がった入力（cos）からは広がった出力}（\(\cos^2\)）が、\textbf{局在入力からは尖った出力}が出る――再生核は入力を歪めず再生するため、入力依存以外の破綻はない。ケース2は
{[}2{]} の図1(a)（cos ベースが全 \(N\) で \(\cos^2\)
に乗る）が数値的裏付けであり、本稿の検証コード §5
にも明示テストとして含める。この\textbf{両ケースの一致}は、機構が入力依存のみで整合する（反証条件をすべてクリアする）ことの最強の証拠である（§5
E）。

\textbf{予想（橋・未証明）}。上の強度プロファイルが、閾値型検出を経て、\textbf{単発クリックの頻度}
\(\propto|\psi_{\rm base}|^2\) として現れる。すなわち
B3（頻度＝強度）を強度プロファイルに帰着し、B1（単発性）を閾値検出に委ねる。この橋は本稿が\textbf{導出しない}部分であり、§6
の系譜 A（Khrennikov）が担う機構に接続する。

\begin{center}\rule{0.5\linewidth}{0.5pt}\end{center}

\subsection{4.
なぜ決定論と確率が両立するか：任意関数の方法}\label{ux306aux305cux6c7aux5b9aux8ad6ux3068ux78baux7387ux304cux4e21ux7acbux3059ux308bux304bux4efbux610fux95a2ux6570ux306eux65b9ux6cd5}

§3 の機構の数学的本体は\textbf{任意関数の方法}（method of arbitrary
functions）である。これは、高速に変化する位相（または力学パラメータ）と、その上の\textbf{任意の滑らかな分布}が、ある極限で\textbf{普遍的な出力分布}へ収束し、入力の詳細によらない、という古典確率論の伝統である
{[}8--10{]}。Poincaré
はルーレット盤の回転を解析し、初期角速度のどんな滑らかな分布も、高速度で各目にほぼ等確率を与えることを示した
{[}8{]}。Diaconis らはコイン投げで同種の結果を与えた
{[}17{]}。この方法の要点は、\textbf{自明でない確率（0でも1でもない）を、より深い層の決定論と両立させる}ことにある。

本稿の構図はこの量子版である。Poincaré
の盤の\textbf{赤黒の交互の弧}（円周上の周期的2値構造）の上で高速回転が等分布するのと、本稿の測定が\textbf{0/180の2値符号}（実軸射影の符号）を読み、高周波ビートが相対位相を等分布させるのとは、同一の数学的対象――周期構造上の高速位相の等分布――である。Feintzeig
ら {[}7{]}
は、ある力学パラメータを初期分布をもつ確率変数として扱うと、長時間かつ古典極限の同時極限で、任意の初期分布に対しボルン則が\textbf{普遍的極限挙動}として出ることを示した。本稿はこの「等分布する力学パラメータ」を、観測器--系の\textbf{相対位相}として物理的に同定する。

したがって、ランダム性は状態にも観測器にも単独で内在しない。それは\textbf{二つの決定論的振動子の相対位相}が測定列の上で等分布することから創発する。これは循環を断つ：どこにも
\(|\psi|^2\)
を仮定せず、可観測性の境界（絶対位相の不在）という独立な事実からゆらぎを出す。

\begin{center}\rule{0.5\linewidth}{0.5pt}\end{center}

\subsection{5.
数値的傍証（検証可能な核）}\label{ux6570ux5024ux7684ux508dux8a3cux691cux8a3cux53efux80fdux306aux6838}

§3 の\textbf{(I) 証明済みの核}を数値で確認する。ベース波
\(\psi_{\rm base}=\cos\varphi+0.5\cos3\varphi-0.3\cos5\varphi\)、\(N=9\)。

\textbf{(A) ビート位相の Weyl 等分布。} 相対位相比を無理数（黄金比
\((\sqrt5-1)/2\)）とし、\(\varphi_0=\)
無理回転で標本化したときの一様分布からの星状不一致（star-discrepancy）：

{\def\LTcaptype{none} % do not increment counter
\begin{longtable}[]{@{}
  >{\raggedright\arraybackslash}p{(\linewidth - 8\tabcolsep) * \real{0.2000}}
  >{\raggedright\arraybackslash}p{(\linewidth - 8\tabcolsep) * \real{0.2000}}
  >{\raggedright\arraybackslash}p{(\linewidth - 8\tabcolsep) * \real{0.2000}}
  >{\raggedright\arraybackslash}p{(\linewidth - 8\tabcolsep) * \real{0.2000}}
  >{\raggedright\arraybackslash}p{(\linewidth - 8\tabcolsep) * \real{0.2000}}@{}}
\toprule\noalign{}
\begin{minipage}[b]{\linewidth}\raggedright
ビート回数 \(N_{\rm beats}\)
\end{minipage} & \begin{minipage}[b]{\linewidth}\raggedright
\(10^2\)
\end{minipage} & \begin{minipage}[b]{\linewidth}\raggedright
\(10^3\)
\end{minipage} & \begin{minipage}[b]{\linewidth}\raggedright
\(10^4\)
\end{minipage} & \begin{minipage}[b]{\linewidth}\raggedright
\(10^5\)
\end{minipage} \\
\midrule\noalign{}
\endhead
\bottomrule\noalign{}
\endlastfoot
star-discrepancy & \(1.4\times10^{-2}\) & \(1.3\times10^{-3}\) &
\(2.6\times10^{-4}\) & \(3.0\times10^{-5}\) \\
\end{longtable}
}

\(N_{\rm beats}\) 増大とともに \(0\)
へ収束（等分布）。一方\textbf{有理比}（\(1/2,\,1/3,\,1/4\)）では不一致が
\(0.25\)--\(0.50\)
で止まり、等分布しない（縞が偏る・反例）。すなわち「観測器が系と一般の（無理比の）振動数をもつ」が等分布の条件である。

\textbf{(B) 再生核の実畳み込み（両ケース）。} v0.2 では
\((S_N*\psi_{\rm base})(\varphi_0)=\int_{-\pi/2}^{\pi/2}S_N(\varphi_0-\varphi)\psi_{\rm base}(\varphi)\,d\varphi\)
を\textbf{実際に数値積分}し、\(\tfrac{\pi}{2}\psi_{\rm base}(\varphi_0)\)
との一致を本稿のコードで自己完結に確認する（{[}2{]}
の再掲に留めない）。三つの代表入力すべてで機械精度一致:

{\def\LTcaptype{none} % do not increment counter
\begin{longtable}[]{@{}
  >{\raggedright\arraybackslash}p{(\linewidth - 4\tabcolsep) * \real{0.3333}}
  >{\raggedright\arraybackslash}p{(\linewidth - 4\tabcolsep) * \real{0.3333}}
  >{\raggedright\arraybackslash}p{(\linewidth - 4\tabcolsep) * \real{0.3333}}@{}}
\toprule\noalign{}
\begin{minipage}[b]{\linewidth}\raggedright
入力 \(\psi_{\rm base}\)
\end{minipage} & \begin{minipage}[b]{\linewidth}\raggedright
\(\max|(S_N*\psi)-\tfrac{\pi}{2}\psi|\)
\end{minipage} & \begin{minipage}[b]{\linewidth}\raggedright
\(\max\big||\,\cdot\,|^2_{\rm norm}-|\psi|^2_{\rm norm}\big|\)
\end{minipage} \\
\midrule\noalign{}
\endhead
\bottomrule\noalign{}
\endlastfoot
ケース1 局在 \(\psi^{\rm loc}=S_N\) & \(1.8\times10^{-15}\) &
\(8.9\times10^{-16}\) \\
ケース1' 多モード \(\cos+0.5\cos3-0.3\cos5\) & \(4.4\times10^{-16}\) &
\(4.4\times10^{-16}\) \\
ケース2 広がり \(\cos\varphi\) & \(4.4\times10^{-16}\) &
\(3.3\times10^{-16}\) \\
\end{longtable}
}

\textbf{同一の再生核}から、出力差は入力の差のみ：局在入力 \(\to\)
尖った出力、広がり入力 \(\to\) 広がった出力（§3
の両ケース導出の数値確認）。

\textbf{バンド条件（\(N\) 依存）}。再生核 \(S_N\) が再生するのは
\(|k|\le N\) のモードに限る。これを \(N\) を振って確認する（各セルは
\(\max|(S_N*\psi)-\tfrac{\pi}{2}\psi|\)）：

{\def\LTcaptype{none} % do not increment counter
\begin{longtable}[]{@{}llllll@{}}
\toprule\noalign{}
入力 & \(N=1\) & \(N=3\) & \(N=5\) & \(N=9\) & \(N=31\) \\
\midrule\noalign{}
\endhead
\bottomrule\noalign{}
\endlastfoot
\(\cos\varphi\)（モード1） & 2e-16 & 2e-16 & 4e-16 & 4e-16 & 2e-16 \\
多モード（1,3,5） & 1.2 & 0.47 & 4e-16 & 4e-16 & 7e-16 \\
\end{longtable}
}

\(\cos\) は\textbf{全 \(N\) で厳密}、多モード（最高次数5）は
\textbf{\(N\ge5\) で厳密・\(N<5\) で取りこぼし}。これは前提
(ii)「観測できない位相差は無視」＝打ち切りが帯域を覆えば厳密、の現れである。なお
(B) は前版の「\((\pi/2)\psi\) を代入した恒等」ではなく、\(S_N\)
を\textbf{実際に数値畳み込みした自己完結の独立検証}である（前版のタウトロジーを解消）。

\textbf{(C) 単一複素モード。}
\(\psi=e^{ik\varphi}\)（\(k=1,3,5\)）では二乗読み出しの位置依存の幅が
\(\sim3\times10^{-15}\)（=0、一様
\(=|e^{ik\varphi}|^2=1\)）。鋭い単一振動数のモードは位置情報をもたず完全非局在、という
\(\Delta\nu\cdot\Delta\varphi\gtrsim1\) の現れ。

\textbf{(D) 役割分担の数値的分離（Weyl だけでは偏りは出ない）。} §3
の役割分担を直接確かめる。無理回転で等分布させた
\(\varphi_0\)（\(2\times10^6\)
点）の\textbf{ヒストグラムそのものは平坦}で、一様密度 \(1/\pi\)
からの最大偏差は
\(1.0\times10^{-5}\)――\textbf{等分布は偏りを作らない}。一方、各
\(\varphi_0\) 訪問を強度 \(|\psi_{\rm base}(\varphi_0)|^2\)
に比例する確率で採択する（閾値検出の橋のモンテカルロ模型）とき、採択クリックのヒストグラムは
\(|\psi_{\rm base}|^2\) に一致する（ケース2 cos：偏差
\(2.7\times10^{-3}\)、ケース1' 多モード：\(6.7\times10^{-3}\)、いずれも
MC 統計誤差）。すなわち \(|\psi|^2\)
の偏りは\textbf{再生核の決定論的読み＋橋}から来るのであって、等分布からではない。

\textbf{重要（地位の明示）}：この D2 のモンテカルロは「強度
\(|\psi_{\rm base}(\varphi_0)|^2\)
に比例してクリックさせる」という\textbf{橋(§3
II)を最初から仮定して実装した整合性デモ}であって、\textbf{橋の導出でも検証でもない}。強度に比例採択すれば出力が
\(|\psi|^2\)
になるのは当然であり、ここで示したのは「橋を仮定すれば全体が無矛盾に閉じる（一様測度
× 強度プロファイル × 橋 ＝ \(|\psi|^2\)
頻度）」という一貫性に過ぎない。したがって本稿の地位分離は崩れない：\textbf{核}＝Weyl（証明済み）×
再生核（証明済み）＝一様測度 × 強度プロファイルの\textbf{形}
\(|\psi_{\rm base}|^2\)（D1 と (B) が裏取り）、\textbf{橋}＝強度 →
単発クリック確率（\textbf{予想}、D2
では仮定として置いただけ・閾値検出＝系譜A に委ねる）。D2
は強度がまだ確率でない（指摘：強度プロファイルは古典的時間平均干渉縞）という地位を\textbf{変えない}。

\textbf{(E) 反証バッテリー（整合性条件）。}
機構が「入力依存のみ」で破綻していないことを、破れたら理論が誤りになる条件で確認する：

\begin{enumerate}
\def\labelenumi{\arabic{enumi}.}
\tightlist
\item
  \textbf{広がり入力 → 尖った出力になっていないか}：cos
  入力の出力ピーク度（最大/平均）\(=1.96\)（低い＝広がったまま）。\textbf{OK}（破綻なし）。
\item
  \textbf{局在入力 → 広がりすぎていないか}：\(S_N\) 入力の出力ピーク度
  \(=9.80\)（高い＝尖ったまま）。\textbf{OK}。
\item
  \textbf{Weyl 単独で二乗の偏りが出ていないか}：(D1)
  で平坦を確認。\textbf{OK}。
\item
  \textbf{ケース2 がケース1 と独立に偽になりうるか}：両者は同一恒等式
  \((S_N*\psi)=\tfrac{\pi}{2}\psi\)
  の特例で、独立に反証されえない（同じ橋の上）。\textbf{OK}。
\end{enumerate}

四条件すべてをクリアする。これは「同一機構・入力依存のみで両ケースが正しく出る」ことの確認であり、\textbf{予想の整合性に対する積極的な（攻めの）証拠}である。

(A)（Weyl＝一様測度）と (B)（再生核＝形の決定論的読み）の積が、§3
の枠付き等式――位置分解強度プロファイルの\textbf{形}
\(=|\psi_{\rm base}|^2\)――を構成する。これが本稿の\textbf{（等分布の仮定
O1 の下で）証明済みの核}である。(D)(E)
は役割分担と整合性を裏付ける。残る橋（強度→単発確率）は数値では確認できない種類の主張であり、§6・§8
で扱う。

\begin{center}\rule{0.5\linewidth}{0.5pt}\end{center}

\subsection{6.
先行研究との関係}\label{ux5148ux884cux7814ux7a76ux3068ux306eux95a2ux4fc2}

本稿の予想は二つの確立した系譜の交点にある。

\subsubsection{6.1 系譜A：Khrennikov の PCSFT ＋
閾値検出（物理的描像）}\label{ux7cfbux8b5cakhrennikov-ux306e-pcsft-ux95beux5024ux691cux51faux7269ux7406ux7684ux63cfux50cf}

Khrennikov の prequantum classical statistical field theory
では、量子系は古典ランダム場の象徴的表現であり、シュレディンガー動力学は古典場の線形動力学の特殊形、測定は古典場と検出器の相互作用、そして「波束の収縮」＝不連続性は\textbf{閾値型検出器の使用という自明な起源}をもつ
{[}6{]}。検出閾値はプリ量子信号の平均エネルギーで決まり、適切に較正された閾値検出器による古典信号の測定としてボルン則が再現される
{[}3,4{]}。背景場（真空ゆらぎ）の存在が鍵要素である {[}4{]}。

本稿との一致：\textbf{離散検出は状態でなく検出器側から来る}（B1）、\textbf{測定器は固有の物理をもつ}。相違：Khrennikov
のランダム性の源は\textbf{確率的なプリ量子場}（\(L^2\) 値 Wiener
過程・ガウス場）であり
{[}5{]}、ボルン則は\textbf{近似的}で破れを予言する。本稿のランダム性の源は\textbf{決定論的ビート}（相対位相の等分布）であり、再生核性により形は\textbf{厳密}である。本稿の予想の「橋」（強度→単発クリック）は、まさに
Khrennikov
型の閾値検出が担う機構に対応する。すなわち本稿は系譜Aを\textbf{機構の供給源}として引く。

\subsubsection{6.2
系譜B：任意関数の方法（数学的機構）}\label{ux7cfbux8b5cbux4efbux610fux95a2ux6570ux306eux65b9ux6cd5ux6570ux5b66ux7684ux6a5fux69cb}

§4 のとおり、本稿の等分布論法は Poincaré {[}8{]}・Hopf {[}9{]}・Engel
{[}10{]} に始まる任意関数の方法そのものであり、その量子適用が Feintzeig
ら {[}7{]}
である。本稿固有の寄与は、彼らが抽象的に置く「確率変数とする力学パラメータ」を、観測器--系の\textbf{ヘテロダイン相対位相}として物理的に同定し（{[}2{]}
が {[}15{]}
のラムゼー／ヘテロダイン計量に接地させた線の延長）、かつ強度プロファイルの形を再生核性で厳密に固定する点にある。

\subsubsection{6.3
隣接系譜と典型性}\label{ux96a3ux63a5ux7cfbux8b5cux3068ux5178ux578bux6027}

基盤側では、't Hooft の決定論的（セルラーオートマトン）量子力学 {[}12{]}
が離散基層に並び、Buniy・Hsu・Zee「離散性と量子力学における確率の起源」{[}11{]}
が離散化第一の立場に連なる。また Zurek の envariance {[}13{]} と Everett
系の典型性は、任意関数の方法と\textbf{同じ「対称性／典型性 →
重み」族}に属する。本稿の等分布（B3）・{[}2{]}
の等振幅（再生核）・envariance
の等確率枝は、一つの族の異なる顔と見なせる。

\subsubsection{6.4
位置づけのまとめ}\label{ux4f4dux7f6eux3065ux3051ux306eux307eux3068ux3081}

本稿の貢献は交点として明確である：\textbf{Khrennikov
型の測定誘起の離散性（系譜A）＋
任意関数の方法による等分布の普遍重み（系譜B）＋
再生核による形の厳密性（{[}2{]}）＋
等分布する高速変数のヘテロダイン同定（本稿）}。この組み合わせは既存研究に見当たらない。

\begin{center}\rule{0.5\linewidth}{0.5pt}\end{center}

\subsection{\texorpdfstring{7. 検証可能な予言：有限 \(N\)
補正＝ボルン則の破れ}{7. 検証可能な予言：有限 N 補正＝ボルン則の破れ}}\label{ux691cux8a3cux53efux80fdux306aux4e88ux8a00ux6709ux9650-n-ux88dcux6b63ux30dcux30ebux30f3ux5247ux306eux7834ux308c}

系譜A・Bがいずれもボルン則を\textbf{極限／近似}でしか出さないことは示唆的である：Khrennikov
は近似（かつ破れを予言）{[}3{]}、任意関数の方法は長時間かつ古典極限の同時極限
{[}7{]}。本稿の予想も、ボルン則を\textbf{等分布極限}（無理ビート・\(N\)
大）での普遍的・典型的挙動として述べるのが妥当であり、\textbf{有限では補正を持つ}。

具体的に、{[}2{]} の連続区間の核 \(S_N\)
を\textbf{有限格子上の奇数倍音核}（離散時空の Nyquist
で硬く切れた核）に置き換えると、(i) 勾配エネルギー \(\sum\nu^2\)
が厳密に有限化し、(ii) 再生がもはや厳密でなく \(O(1/N)\)
の補正を持つ。ここで \(N=R/\ell\) はスケール比（系の局在幅 \(\ell\)
と区間
\(R\)）である。この補正は\textbf{ボルン則からの位置依存の逸脱}として観測可能で、その大きさは
\(\sim1/N=\ell/R\)、最大となるのは帯端（高次モード近傍）と予言される。これは系譜A（Khrennikov
の予言する破れ）と同じ実験的射程にあり、本予想を\textbf{反証可能}にする。

\begin{center}\rule{0.5\linewidth}{0.5pt}\end{center}

\subsection{8. 未解決問題}\label{ux672aux89e3ux6c7aux554fux984c}

本稿が\textbf{導出しない}ものを明示する。

\begin{itemize}
\tightlist
\item
  \textbf{O1（等分布の仮定）}：位置分解強度が \(|\psi|^2\)
  に一致するには、相対位相が無理比で等分布する必要がある（§5
  で確認した条件）。これは自然だが仮定であり、B3
  は完全には消えず「観測器位相の等分布」へ\textbf{移設}される。ただしこれは
  \(|\psi|^2\) を直接仮定するよりはるかに弱い。
\item
  \textbf{O2（B1・単発性）}：ビートは \(\varphi_0\)
  の\textbf{分布}を与えるが、なぜ一回の試行で\textbf{単一}のクリックが立つかは与えない。これは本稿の射程外の\textbf{別個の機構}であり、系譜A（閾値検出）が担う部分でもある。なお
  §3--§4 の等分布論法（無理回転 →
  Weyl）は\textbf{離散的な測定イベント列}を前提とする（連続時間では無理性なしに掃引が一様化し、有理比の反例も消える）。\textbf{本稿はこの離散測定列を仮定として残し、その起源（離散化の機構）には立ち入らない。}ゆえに
  B1 の離散クロックは B2・B3 の機構の\textbf{前提}でもあり、O2
  は他の未解決問題から完全には独立でない。等分布（O1）と単発性（O2）は、いずれも「離散測定イベントの存在」という共通の上流を分け持つ。
\item
  \textbf{O3（二乗のべき）}：振幅を確率へ変える\textbf{指数2}は、複素構造（実2次元）の面積測度（Robertson
  的）から来るべきもので、ビートからは出ない。{[}2{]} の複素モジュラス
  \(Z\bar Z=\mathrm{Re}^2+\mathrm{Im}^2\) の議論に属する。
\item
  \textbf{O4（スケール帰属）}：\(\sum\nu_n^2=\nu^2\)
  型の写像において、\(\nu\) と内部 \(\nu_n\)
  のどちらが上位スケールかを確定する必要がある。観測器が見る単一振動数
  \(\nu\) と内部微細位相の遮蔽の向きは、この帰属に依存する。
\end{itemize}

\begin{center}\rule{0.5\linewidth}{0.5pt}\end{center}

\subsection{9. 結論}\label{ux7d50ux8ad6}

本稿は、局在核論文 {[}2{]}
が公準として残した\textbf{ランダム性の起源}に対する機構の候補を\textbf{予想}として提示した。中心の観察は、走査位相
\(\varphi_0\)
が観測器の位相であり、観測器が有限振動数をもつ物理対象ゆえ絶対参照が存在しない（主観空間の前提）こと、ゆえに測定列が観測器--系の相対位相（ビート）を走査し、それが無理比で等分布する（Weyl）ことである。

論理は二段である。\textbf{（O1
の下で）証明済みの核}：再生核性（{[}2{]}）と Weyl
等分布から、ビートが掃く\textbf{位置分解強度プロファイル}が
\(|\psi_{\rm base}|^2\) に一致する（等分布 O1 を認めれば追加仮定なく、§5
で数値確認）。\textbf{予想の橋}：この強度が閾値検出を経て単発クリックの頻度
\(\propto|\psi_{\rm base}|^2\)
として現れる（本稿は導出しない、系譜Aに接続）。

この予想は二つの確立した系譜――Khrennikov の PCSFT ＋
閾値検出（物理的描像）と任意関数の方法（数学的機構）――の交点にあり、本稿固有の寄与は、等分布する高速変数を観測器--系のヘテロダイン・ビートとして物理同定し、再生核性で形を厳密化する点にある。さらに有限
\(N\)（離散格子）での \(O(1/N)\sim\ell/R\)
補正が、ボルン則からの測定可能な逸脱という反証可能な予言を与える。

残るのは
B1（単発性）・指数2・スケール帰属であり、本稿はこれらを未解決として正直に切り出す。ボルン則を「導出した」のではなく、\textbf{その統計的起源を、観測器の有限性と相対位相の等分布という二つの独立に正しい要素へ移設する候補}を、検証可能な形で提示した、というのが本稿の射程である。

\begin{center}\rule{0.5\linewidth}{0.5pt}\end{center}

\subsection{付録A：最小検証の素描}\label{ux4ed8ux9332aux6700ux5c0fux691cux8a3cux306eux7d20ux63cf}

ベース波
\(\psi_{\rm base}=\cos\varphi+0.5\cos3\varphi-0.3\cos5\varphi\)、\(N=9\)。

\begin{enumerate}
\def\labelenumi{\arabic{enumi}.}
\tightlist
\item
  \textbf{Weyl 等分布}：相対位相比 \(r=(\sqrt5-1)/2\)（無理）で
  \(\varphi_0^{(n)}=\big(r\,n \bmod 1\big)\cdot\pi-\pi/2\)
  を生成し、星状不一致が \(N_{\rm beats}\to\infty\) で \(0\)
  に収束することを確認（§5 表）。有理 \(r\) では収束しない（反例）。
\item
  \textbf{位置分解強度}：各 \(\varphi_0\) で
  \(|(\tfrac{\pi}{2})\psi_{\rm base}(\varphi_0)|^2\)
  を規格化し、\(|\psi_{\rm base}|^2\) と機械精度一致を確認。
\item
  \textbf{単一複素モード}：\(\psi=e^{ik\varphi}\)
  で二乗読み出しが位置一様（\(=1\)）を確認。
\end{enumerate}

これらは検証コードで機械精度／古典定理として再現する。橋（強度→単発確率）は数値検証の対象外であり、概念的・実験的（有限
\(N\) 補正）にのみ問える。

\begin{center}\rule{0.5\linewidth}{0.5pt}\end{center}

\subsection{参考文献}\label{ux53c2ux8003ux6587ux732e}

{[}1{]} M. Born, ``Zur Quantenmechanik der Stoßvorgänge,'' \emph{Z.
Phys.} \textbf{37}, 863 (1926). {[}2{]} 木原 範昭,
「局在奇数倍音波の再生核性によるボルン分布の形の導出 ―
残す公準を二乗則とランダム性に絞る」, Zenodo, v0.1 (2026-06-27), Concept
DOI: 10.5281/zenodo.20965526 / Version DOI:
10.5281/zenodo.20965527（本稿の出発点：再生核による形の厳密導出）。
{[}3{]} A. Khrennikov, ``Born's rule from measurements of classical
signals by threshold detectors which are properly calibrated,''
arXiv:1105.4269 (2011). {[}4{]} A. Khrennikov, ``Quantum Probabilities
and Violation of CHSH-Inequality from Classical Random Signals and
Threshold Type Detection Scheme,'' \emph{Prog. Theor. Phys.}
\textbf{128}, 31 (2012). {[}5{]} A. Khrennikov, ``Born's rule from
statistical mechanics of classical fields: from hitting times to quantum
probabilities,'' arXiv:1212.0756 (2012). {[}6{]} A. Khrennikov,
``Measurement problem: from de Broglie to theory of classical random
fields interacting with threshold detectors,'' arXiv:1301.0011 (2013).
{[}7{]} L. Bonds, B. Burson, K. Cicchella, B. H. Feintzeig, Lynnx, A.
Yusaini, ``Quantum Probability via the Method of Arbitrary Functions,''
arXiv:2409.16457 (2024).（著者「Lynnx」は arXiv 登録どおりの単名表記。）
{[}8{]} H. Poincaré, \emph{Calcul des probabilités} (Gauthier-Villars,
Paris, 1896)（méthode des fonctions arbitraires：ルーレット盤）。
{[}9{]} E. Hopf, ``On causality, statistics and probability,'' \emph{J.
Math. Phys.} \textbf{13}, 51 (1934). {[}10{]} E. Engel, \emph{A Road to
Randomness in Physical Systems}, Lecture Notes in Statistics \textbf{71}
(Springer, 1992). {[}11{]} R. V. Buniy, S. D. Hsu, A. Zee,
``Discreteness and the origin of probability in quantum mechanics,''
\emph{Phys. Lett. B} \textbf{640}, 219 (2006). {[}12{]} G. 't Hooft,
\emph{The Cellular Automaton Interpretation of Quantum Mechanics},
Fundamental Theories of Physics \textbf{185} (Springer, 2016). {[}13{]}
W. H. Zurek, ``Probabilities from entanglement, Born's rule
\(p_k=|\psi_k|^2\) from envariance,'' \emph{Phys. Rev.~A} \textbf{71},
052105 (2005). {[}14{]} A. Valentini and H. Westman, ``Dynamical origin
of quantum probabilities,'' \emph{Proc. R. Soc. A} \textbf{461}, 253
(2005). {[}15{]} N. F. Ramsey, ``A Molecular Beam Resonance Method with
Separated Oscillating Fields,'' \emph{Phys. Rev.} \textbf{78}, 695
(1950). {[}16{]} H. Weyl, ``Über die Gleichverteilung von Zahlen mod.
Eins,'' \emph{Math. Ann.} \textbf{77}, 313 (1916)（一様分布定理）。
{[}17{]} P. Diaconis, S. Holmes, R. Montgomery, ``Dynamical bias in the
coin toss,'' \emph{SIAM Rev.} \textbf{49}, 211 (2007).

\begin{center}\rule{0.5\linewidth}{0.5pt}\end{center}

\emph{(本稿は査読対話ノート（木原範昭 × アイリス,
2026-06-27）に由来する予想論文である。（O1
の下で）証明済みの核（再生核性 ＋ Weyl 等分布による位置分解強度
\(=|\psi_{\rm base}|^2\)）と、明示的に区切った未証明の橋（強度 →
単発クリック確率）からなり、後者は系譜A（Khrennikov
型閾値検出）に接続する。B1・指数2・スケール帰属は未解決として残す。)}

\end{document}
